\documentclass[12pt,a4paper]{article}

\usepackage[utf8]{inputenc}
\usepackage[T1]{fontenc}
\usepackage[brazilian]{babel}
\usepackage{geometry}
\usepackage{booktabs}
\usepackage{array}
\usepackage{tabularx}
\usepackage{xcolor}
\usepackage{titlesec}
\usepackage{fancyhdr}
\usepackage{setspace}
\usepackage{parskip}
\usepackage{hyperref}
\usepackage{indentfirst}
\usepackage{amsmath}
\usepackage{enumitem}
\usepackage{colortbl}

\setlength{\headheight}{25pt}
\setlength{\emergencystretch}{3em}

\geometry{
  a4paper,
  top=3cm,
  bottom=2cm,
  left=3cm,
  right=2cm
}

\definecolor{azul}{RGB}{31, 92, 153}
\definecolor{azulclaro}{RGB}{213, 232, 240}
\definecolor{cinza}{RGB}{85, 85, 85}

\titleformat{\section}
  {\large\bfseries\color{azul}}
  {\thesection.}{0.5em}{}
  [\vspace{2pt}{\color{azul}\titlerule[1pt]}\vspace{4pt}]

\titleformat{\subsection}
  {\normalsize\bfseries\color{azul!80!black}}
  {\thesubsection}{0.5em}{}

\titlespacing*{\section}{0pt}{18pt}{8pt}
\titlespacing*{\subsection}{0pt}{12pt}{4pt}

\pagestyle{fancy}
\fancyhf{}
\fancyhead[L]{\small\color{cinza} Projeto Aplicado em Ciência e Inovação}
\fancyhead[R]{\small\color{cinza} \thepage}
\renewcommand{\headrulewidth}{0.4pt}
\renewcommand{\headrule}{\hbox to\headwidth{\color{azul}\leaders\hrule height \headrulewidth\hfill}}

\onehalfspacing

\hypersetup{
  colorlinks=true,
  linkcolor=azul,
  urlcolor=azul,
  citecolor=azul
}

\begin{document}

% CAPA
\thispagestyle{empty}
\begin{center}
  \vspace*{1cm}
  {\large\scshape\color{cinza} Universidade Federal do Rio Grande do Sul}\\[4pt]
  {\normalsize\color{cinza} Ciência da Computação --- Tópicos Especiais}\\[2cm]

  {\color{azul}\rule{\linewidth}{2pt}}\\[0.6cm]
  {\LARGE\bfseries\color{azul}
    Análise de Vulnerabilidade e Simulação de\\[6pt]
    Falhas em Cascata na Rede Aeroportuária\\[6pt]
    Brasileira: Uma Abordagem por\\[6pt]
    Teoria dos Grafos
  }\\[0.6cm]
  {\color{azul}\rule{\linewidth}{2pt}}\\[1cm]

  {\normalsize\color{cinza} Trabalho final --- relatório técnico-científico}\\[2cm]

  {\large Tobias Cadoná Marion}\\[4pt]
  {\large Vitor Godoy de Fraga}\\[4pt]
  {\large Juan Marcelo Trabaina}\\[3cm]

  {\normalsize\color{cinza} Porto Alegre, 2026}
\end{center}

\newpage
\tableofcontents
\newpage

% 1. RESUMO
\section{Resumo e Palavras-Chave}

A rede aeroportuária brasileira constitui uma infraestrutura estratégica para a integração
territorial, econômica e social do país. Em razão da extensão continental do Brasil e da
concentração de conexões em poucos aeroportos de grande porte, interrupções operacionais em
hubs relevantes podem produzir efeitos sistêmicos sobre a conectividade da malha aérea. Este
trabalho desenvolveu um \textbf{artefato computacional aberto e reproduzível} que identifica os
aeroportos estruturalmente mais críticos da malha aérea brasileira e simula o impacto de sua
remoção. A partir de microdados públicos da ANAC referentes aos doze meses de 2025
(1{,}43 milhão de etapas de voo), a malha foi representada como um grafo ponderado com
155 aeroportos e 622 rotas, sobre o qual se calcularam indicadores de centralidade, um ranking
de criticidade e duas famílias de simulação: (i) a remoção topológica de nós, comparando falhas
aleatórias e ataques direcionados; e (ii) um modelo exploratório de realocação de demanda sob
restrição de capacidade. Os resultados confirmam o comportamento esperado de uma rede
\textit{scale-free}: o aeroporto de Viracopos (VCP) emerge como o nó mais crítico --- em
concordância com a literatura independente --- e o ataque direcionado fragmenta a malha de forma
muito mais acentuada que a falha aleatória de mesma magnitude. O ranking dos dez aeroportos mais
críticos mostrou-se robusto a perturbações nos pesos do índice (VCP sempre em primeiro lugar e ao
menos 9 dos 10 aeroportos retidos). No modelo
de capacidade, a remoção dos dez maiores hubs deixa de 29\% a 61\% da demanda de passageiros sem
realocação viável, conforme a folga de capacidade assumida. Diferentemente das análises
acadêmicas pontuais já publicadas sobre a rede brasileira e das ferramentas comerciais
proprietárias, a contribuição central é a entrega de um artefato aberto, testado e auditável ---
uma \textbf{inovação de processo} --- que transforma dados públicos dispersos em indicadores
objetivos e reproduzíveis de resiliência da malha aérea nacional.

\medskip
\noindent\textbf{Palavras-chave:} Redes Complexas; Teoria dos Grafos; Falhas em Cascata;
Rede Aeroportuária; Scale-Free; Simulação Computacional; ANAC.

% 2. INTRODUÇÃO
\section{Introdução e Contextualização}

\subsection{Contexto}

A aviação civil brasileira movimenta milhões de passageiros por ano e exerce papel essencial na
ligação entre centros urbanos, regiões produtivas e localidades geograficamente isoladas. A malha
aérea nacional funciona como um sistema interdependente: a operação de um aeroporto influencia
diretamente a circulação de passageiros, a distribuição de rotas e a capacidade de atendimento de
outros aeroportos conectados.

Nesse tipo de infraestrutura, alguns aeroportos concentram volume elevado de conexões e
passageiros, assumindo função de hubs. Sob a perspectiva da teoria dos grafos e das redes
complexas, essa característica torna a rede relativamente resistente a falhas pontuais em nós
periféricos, mas vulnerável à indisponibilidade de nós centrais. Assim, o fechamento de aeroportos
como Guarulhos (GRU), Congonhas (CGH), Brasília (BSB) ou Viracopos (VCP) pode gerar efeitos que
ultrapassam o local da falha inicial, afetando rotas, conexões e aeroportos alternativos.

\subsection{O Problema}

O problema investigado neste projeto consiste em identificar quais aeroportos exercem maior
influência estrutural sobre a rede aeroportuária brasileira e estimar os efeitos de sua
indisponibilidade sobre a conectividade da malha aérea. Embora existam dados públicos sobre
movimentação aérea, ainda são limitadas as ferramentas abertas voltadas à simulação de falhas e
à análise de vulnerabilidade da rede aeroportuária nacional.

Dessa forma, a questão orientadora do estudo é: quais aeroportos brasileiros apresentam maior
criticidade estrutural e qual o impacto de sua remoção hipotética sobre a conectividade da malha
aérea, comparando-se cenários de falha aleatória e de ataque direcionado?

% 3. JUSTIFICATIVA
\section{Justificativa}

\subsection{Impacto Prático}

O trabalho apresenta relevância prática por aplicar métodos computacionais à análise de uma
infraestrutura crítica nacional. A identificação de aeroportos estruturalmente relevantes pode
apoiar estudos de planejamento, contingência e priorização de investimentos em capacidade
operacional. Além disso, a simulação de cenários de falha permite observar, ainda que em nível
experimental, como a indisponibilidade de determinados aeroportos pode afetar regiões
dependentes do transporte aéreo.

Embora a estrutura da rede aérea brasileira já tenha sido objeto de estudos acadêmicos pontuais
(ver Seção~\ref{sec:lit}) e existam ferramentas comerciais proprietárias de análise do setor,
persiste uma lacuna: a ausência de uma ferramenta \textbf{aberta e reproduzível} que converta os
microdados públicos da ANAC em indicadores de vulnerabilidade auditáveis. É essa lacuna que o
trabalho busca ocupar; a natureza dessa contribuição e seu impacto são discutidos, à luz dos
resultados, na Seção~\ref{sec:discussao}.

\subsection{Pertinência da Computação}

A Computação é pertinente ao problema porque a rede aeroportuária envolve múltiplos aeroportos,
rotas, volumes de passageiros e cenários de interrupção. A análise manual dessas relações tende a
ser limitada, especialmente quando se deseja avaliar o comportamento global da rede após a
remoção de um ou mais nós. A modelagem computacional permite representar formalmente essas
relações, calcular métricas estruturais e executar simulações repetíveis, fornecendo resultados
mais consistentes para comparação entre cenários.

% 4. OBJETIVOS
\section{Objetivos}

\subsection{Objetivo Geral}

Desenvolver um protótipo aberto e reproduzível que identifique os aeroportos estruturalmente mais
críticos da malha aérea brasileira, a partir de microdados públicos da ANAC, e simule o impacto de
sua remoção sobre a conectividade da rede, comparando cenários de falha aleatória e de ataque
direcionado.

\subsection{Objetivos Específicos}

\begin{enumerate}[leftmargin=*, label=\textbf{OE\arabic*.}]
  \item Entregar uma base estruturada de rotas e movimentação de passageiros da rede
        aeroportuária brasileira.
  \item Entregar uma representação computacional da malha aérea nacional em formato de grafo
        ponderado.
  \item Produzir indicadores de centralidade e criticidade para os principais aeroportos da rede.
  \item Disponibilizar um simulador experimental de falhas aleatórias e ataques direcionados
        sobre aeroportos selecionados.
  \item Apresentar um relatório técnico com os cenários simulados, os impactos observados e a
        interpretação dos resultados obtidos.
\end{enumerate}

% 5. REVISÃO DE LITERATURA
\section{Revisão de Literatura e Soluções de Mercado}
\label{sec:lit}

\subsection{Fundamentação Teórica}

\textbf{Albert, Jeong e Barabasi (2000)} estabeleceram uma das principais bases teóricas para o
estudo da robustez em redes complexas. Os autores demonstraram que redes \textit{scale-free},
caracterizadas por distribuição desigual de conexões e concentração de ligações em poucos hubs,
tendem a resistir a falhas aleatórias, mas apresentam fragilidade quando os nós mais conectados
são removidos de forma direcionada. Essa discussão fundamenta a análise da criticidade dos
aeroportos centrais na malha aérea brasileira.

\textbf{Correia et al.} desenvolveram um modelo de previsão de demanda aeroportuária baseado
na equação de Huff, expressa por:

\begin{equation}
  E_{ij} = P_{ij} \cdot E_i = \frac{\dfrac{S_j}{T_{ij}^{\,a}}}{\displaystyle\sum_{j}
  \dfrac{S_j}{T_{ij}^{\,a}}} \cdot C_i
\end{equation}

\noindent onde $C_i$ é a demanda total do centro populacional $i$, $S_j$ é o poder de atração
do aeroporto $j$, $T_{ij}$ é o tempo de viagem entre $i$ e $j$, e $a$ é um parâmetro estimado
empiricamente. Essa abordagem motiva o módulo exploratório de redistribuição de demanda do
protótipo, que realoca a demanda dos aeroportos fechados para alternativas próximas com capacidade
disponível; a adoção do modelo de Huff completo, com tempos de viagem, permanece como refinamento
futuro (ver Seção~\ref{sec:limitacoes}).

\textbf{Cumelles, Lordan e Sallan (2021)} analisam falhas em cascata em redes aeroportuárias e
destacam que esses sistemas apresentam características próprias, pois os voos são eventos
programados e não fluxos contínuos. O trabalho propõe uma forma de avaliar o fechamento de
aeroportos e a realocação de voos para alternativas próximas com capacidade disponível. Essa
contribuição orienta a definição dos cenários de simulação adotados neste projeto.

\subsection{Estado da Técnica e Mercado}

No meio acadêmico, a estrutura da rede aérea brasileira já foi estudada. \textbf{Couto et al.
(2015)} caracterizaram a malha nacional como uma rede \textit{scale-free} com propriedades de
pequeno mundo, identificando o aeroporto de Viracopos como o nó mais central da rede de voos
nacionais e mostrando que sua indisponibilidade fragmenta a malha em diversas sub-redes — evidência
de baixa resiliência a ataques direcionados. Esse e outros estudos, contudo, são artigos de
análise: descrevem resultados, mas não disponibilizam uma ferramenta aberta e reproduzível que
outros pesquisadores possam executar e estender.

No plano comercial, ferramentas como OAG Traffic Analyser e Cirium Flightglobal oferecem análises de
mercado, tráfego e planejamento para o setor aéreo, mas são soluções proprietárias e fechadas. A
ANAC disponibiliza bases públicas de movimentação aérea; entretanto, a conversão desses dados em
uma ferramenta experimental, aberta e reproduzível de análise de vulnerabilidade ainda representa
uma lacuna — exatamente o espaço que esta proposta busca ocupar.

% 6. METODOLOGIA
\section{Metodologia}

\subsection{Fonte de Dados}

O estudo utilizou dados públicos disponíveis no portal de dados abertos da ANAC
(\url{https://dados.anac.gov.br}), com foco em registros de voos regulares, pares
origem-destino e movimentação de passageiros. Foram processados os \textbf{doze} arquivos mensais
combinados de 2025, totalizando 2.350.330 linhas brutas, das quais 1.430.477 etapas de voo foram
utilizadas após filtragem (descartadas as etapas internacionais, não regulares, de carga ou sem
passageiros). Capacidades
reais de pista e \textit{slots} não estão disponíveis em base pública consolidada e, por isso,
foram tratadas como parâmetro do modelo (Seção~\ref{sec:capacidade}), não como dado de entrada.

\subsection{Arquitetura e Algoritmia}

A rede aeroportuária foi representada como um grafo ponderado $G = (V, E)$, no qual $V$
corresponde ao conjunto de aeroportos e $E$ ao conjunto de rotas entre eles. O peso das arestas
foi definido de acordo com o volume de passageiros transportados no período de referência. A
construção do grafo e o cálculo das métricas estruturais foram realizados em Python, com os
algoritmos de rede (busca em componentes conexos e \textit{betweenness centrality}) implementados
diretamente sobre a biblioteca padrão da linguagem, sem dependência de bibliotecas externas de
análise de redes.

A criticidade de cada aeroporto foi estimada por um \textbf{índice composto} que combina três
dimensões complementares: a \textit{betweenness centrality} (papel de intermediação nas rotas), a
\textit{degree centrality} (número de conexões diretas) e o volume de passageiros (relevância
operacional). Formalmente,

\begin{equation}
  \text{score}_i = \frac{\text{BC}_i + \text{DC}_i + \widehat{V}_i}{3},
\end{equation}

\noindent onde $\text{BC}_i$ é a \textit{betweenness} (não ponderada, medindo intermediação
estrutural), $\text{DC}_i$ é o grau normalizado e $\widehat{V}_i$ é o volume de passageiros
normalizado por logaritmo. Adotou-se a \textbf{média simples} das três dimensões: na ausência de
uma base empírica para privilegiar uma métrica sobre as demais, o peso igual é a escolha neutra,
e evita que a criticidade dependa de uma única dimensão. A \textbf{calibração dos pesos} ---
atribuir maior ou menor importância a cada dimensão a partir de dados ou de critério especialista
--- permanece como \textbf{trabalho futuro}; a robustez do ranking a variações nos pesos é
avaliada por análise de sensibilidade (Seção~\ref{sec:resultados}).

O simulador implementou dois cenários de remoção de nós:

\begin{itemize}[leftmargin=*]
  \item \textbf{Falha aleatória:} remoção de aeroportos selecionados uniformemente ao acaso, com
        semente fixa para garantir reprodutibilidade.
  \item \textbf{Ataque direcionado:} remoção dos aeroportos em ordem decrescente do índice de
        criticidade.
\end{itemize}

A cada remoção, foram calculados indicadores de conectividade da rede, como o número de componentes
conectados, o tamanho relativo do maior componente conectado, a quantidade de aeroportos isolados e
a fração do volume de passageiros ainda atendida pelas rotas remanescentes. A comparação entre os
dois cenários permite observar a assimetria entre falhas aleatórias e ataques direcionados,
característica esperada de redes \textit{scale-free}.

\subsection{Modelo de Realocação de Demanda sob Capacidade}
\label{sec:capacidade}

Como camada adicional à análise topológica, implementou-se um \textbf{modelo exploratório de
realocação de demanda sob restrição de capacidade}, adaptado do princípio de Cumelles, Lordan e
Sallan (2021): ao fechar um aeroporto, sua demanda de passageiros é redistribuída para as
alternativas próximas com capacidade disponível, e não propagada continuamente pelas conexões da
rede. A carga de cada aeroporto $i$ é a sua demanda observada $L_i$ (embarcados mais
desembarcados), e a capacidade é modelada como

\begin{equation}
  C_i = (1 + \alpha)\,L_i,
\end{equation}

\noindent onde $\alpha \ge 0$ é a \textbf{folga operacional} --- a fração de demanda adicional que
o aeroporto absorve acima de sua operação normal. Removido um aeroporto, sua demanda é alocada nas
alternativas vivas (por padrão, as geograficamente mais próximas), preenchendo a capacidade
residual de cada candidato até esgotá-la; a demanda que não encontra capacidade disponível é
contabilizada como \textbf{demanda não realocada} (\textit{stranded}). Por construção, nenhum
aeroporto recebe carga acima de $C_i$ --- não há propagação de sobrecarga. O comportamento do
modelo foi examinado por uma \textbf{varredura de $\alpha$}, adotando-se $\alpha = 0{,}20$ (folga
de 20\%) como cenário de referência (Seção~\ref{sec:resultados} e Seção~\ref{sec:limitacoes}).

\subsection{Ambiente de Desenvolvimento}

\begin{itemize}[leftmargin=*]
  \item \textbf{Linguagem:} Python 3.11+ (biblioteca padrão, sem dependências externas)
  \item \textbf{Processamento:} módulos \texttt{csv}, \texttt{json} e \texttt{collections} para
        leitura dos microdados e construção do grafo
  \item \textbf{Visualização:} mapa interativo em HTML autocontido com D3.js
  \item \textbf{Interface:} aplicação de linha de comando (CLI) para consulta do ranking e das
        simulações
  \item \textbf{Controle de versão:} Git/GitHub
\end{itemize}

\subsection{Critérios de Validação}

Adotou-se como critério de validação a comparação entre o comportamento da ferramenta e as
propriedades esperadas para redes \textit{scale-free}, documentadas na literatura. Em particular,
verificou-se se: (i) hubs centrais como Viracopos (VCP), Guarulhos (GRU) e Brasília (BSB) emergem
no topo do ranking de criticidade; (ii) o ataque direcionado fragmenta a rede de forma muito mais
acentuada que a falha aleatória de mesma magnitude. Esse é o comportamento esperado segundo Couto
et al. (2015), que reportam a rede aérea brasileira como não-resiliente a ataques direcionados e
apontam Viracopos como nó crítico cuja remoção fragmenta a malha. Ambos os critérios foram
satisfeitos pelos resultados (Seção~\ref{sec:resultados}): a coerência entre o comportamento da
ferramenta e esses achados independentes constitui evidência de validade.

\subsection{Verificação e Confiabilidade}

É preciso distinguir \textbf{verificação} (o programa faz o que a especificação define) de
\textbf{validação} (o modelo corresponde à realidade). A confiabilidade da ferramenta apoia-se em
três pilares de verificação, deixando a validação empírica como trabalho futuro.

\textbf{(i) Verificação automatizada.} O modelo de realocação sob capacidade é coberto por um
conjunto de testes automatizados que asseguram sua consistência interna: o caso-limite de
capacidade infinita ($\alpha \to \infty$) reproduz exatamente a simulação topológica de ataque
(teste discriminador); a redução da folga $\alpha$ aumenta monotonicamente a demanda não realocada;
e nenhum aeroporto recebe carga acima de sua capacidade. Registra-se ainda, como teste, o
\textit{resultado negativo} que motivou o desenho: a propagação ingênua de sobrecarga colapsa a
rede de forma praticamente insensível a $\alpha$, comportamento coerente com a observação de
Cumelles et al. (2021) de que modelos de fluxo contínuo são inadequados a redes aéreas.

\textbf{(ii) Coerência com a literatura.} Os resultados são confrontados com achados independentes
(Albert et al., 2000; Couto et al., 2015), conforme o Plano de Validação.

\textbf{(iii) Reprodutibilidade.} Todos os indicadores e tabelas são regeneráveis a partir dos
dados públicos e dos \textit{scripts} versionados, sem dependências externas, o que permite
auditoria e replicação independentes.

A confiabilidade aqui afirmada é, portanto, de \emph{consistência interna, coerência com a
literatura e reprodutibilidade} --- não de validação empírica contra dados operacionais reais, que
depende de capacidades declaradas de pista/\textit{slots} e da calibração de $\alpha$, e permanece
como trabalho futuro.

% 7. RESULTADOS
\section{Resultados}
\label{sec:resultados}

\subsection{Caracterização da Rede}

A partir de 1.430.477 etapas de voo úteis (extraídas de 2.350.330 registros brutos), o grafo
resultante reúne \textbf{155 aeroportos} e \textbf{622 rotas} não direcionais, com volume total
de \textbf{100.873.591 passageiros} associado às arestas. A rede é \textbf{conexa} (um único
componente conectado abrangendo os 155 aeroportos), confirmando a ausência de aeroportos
isolados na operação regular de 2025. Essa estrutura --- poucos nós de grau muito elevado e uma
longa cauda de aeroportos pouco conectados --- é compatível com a caracterização da malha aérea
brasileira como rede \textit{scale-free} por Couto et al. (2015).

\subsection{Ranking de Criticidade e sua Robustez}

A Tabela~\ref{tab:ranking} apresenta os dez aeroportos mais críticos segundo o índice composto.
O aeroporto de \textbf{Viracopos (VCP)} ocupa a primeira posição, com a maior \textit{betweenness}
da rede (0{,}286), seguido por Confins (CNF) e Guarulhos (GRU). Vale notar que GRU, apesar de
concentrar o maior volume de passageiros (29{,}9 milhões), aparece em terceiro: a criticidade
\textit{estrutural} não coincide com o mero volume, e sim com o papel de intermediação ---
exatamente o que o índice composto pretende capturar.

\begin{table}[h!]
\centering
\renewcommand{\arraystretch}{1.3}
\small
\begin{tabular}{clrrrr}
\toprule
\textbf{\#} & \textbf{Aeroporto} & \textbf{Grau} & \textbf{Passageiros} & \textbf{Betweenness} & \textbf{Score} \\
\midrule
1  & VCP --- Campinas/SP       & 79 & 11.717.415 & 0{,}2859 & 0{,}5815 \\
2  & CNF --- Confins/MG        & 64 & 12.216.388 & 0{,}2162 & 0{,}5266 \\
3  & GRU --- Guarulhos/SP      & 62 & 29.909.827 & 0{,}1082 & 0{,}5036 \\
4  & REC --- Recife/PE         & 47 &  9.160.522 & 0{,}1263 & 0{,}4542 \\
5  & CGH --- São Paulo/SP      & 48 & 24.493.686 & 0{,}0420 & 0{,}4474 \\
6  & BSB --- Brasília/DF       & 41 & 15.626.199 & 0{,}0468 & 0{,}4251 \\
7  & MAO --- Manaus/AM         & 32 &  2.909.349 & 0{,}1896 & 0{,}4207 \\
8  & GIG --- Rio de Janeiro/RJ & 38 & 11.572.338 & 0{,}0311 & 0{,}4076 \\
9  & SSA --- Salvador/BA       & 33 &  7.128.371 & 0{,}0548 & 0{,}3953 \\
10 & FOR --- Fortaleza/CE      & 29 &  5.490.756 & 0{,}0596 & 0{,}3831 \\
\bottomrule
\end{tabular}
\caption{Dez aeroportos mais críticos por índice composto de criticidade.}
\label{tab:ranking}
\end{table}

O índice usa \textbf{pesos iguais} (média simples), por ser a escolha neutra na ausência de
calibração. Para verificar se o resultado depende dessa escolha, o top-10 foi recomputado sob
\textbf{análise de sensibilidade}, perturbando-se os pesos para enfatizar cada uma das três
dimensões e testando-se também ponderações alternativas. \textbf{VCP permaneceu como nó mais
crítico em todas as configurações}, e o top-10 manteve-se estável: em metade das perturbações o
conjunto foi idêntico ao base (interseção 10/10, Jaccard $1{,}00$) e, no restante, 9 dos 10
aeroportos foram retidos (Jaccard $0{,}82$), com a troca de um único aeroporto na fronteira do
ranking (entre Manaus e Fortaleza, conforme se reforce ou não a \textit{betweenness}). O
resultado, portanto, \textbf{não é dirigido pela escolha exata dos pesos}: os aeroportos
estruturalmente centrais emergem sob qualquer ponderação razoável das três dimensões.

\subsection{Ataque Direcionado \textit{versus} Falha Aleatória}

A Tabela~\ref{tab:ataque} compara, no décimo passo de remoção, o ataque direcionado (remoção dos
dez aeroportos mais críticos) e a falha aleatória (sorteio com semente fixa). A assimetria é
expressiva: o ataque direcionado fragmenta a rede em 48 componentes e isola 37 aeroportos,
reduzindo o maior componente a 87 dos 155 aeroportos (56{,}1\%); a falha aleatória de mesma
magnitude preserva um maior componente de 132 aeroportos (85{,}2\%), com apenas 10 componentes e
5 isolados. A primeira fragmentação ocorreu já no \textbf{passo 1} do ataque (remoção de VCP),
contra o \textbf{passo 6} no cenário aleatório.

\begin{table}[h!]
\centering
\renewcommand{\arraystretch}{1.3}
\small
\begin{tabular}{lccc}
\toprule
\textbf{Cenário (passo 10)} & \textbf{Maior componente} & \textbf{Componentes} & \textbf{Isolados} \\
\midrule
Ataque direcionado & 87 (56{,}1\%) & 48 & 37 \\
Falha aleatória    & 132 (85{,}2\%) & 10 & 5 \\
\bottomrule
\end{tabular}
\caption{Degradação da conectividade após dez remoções, por cenário.}
\label{tab:ataque}
\end{table}

Esse é exatamente o comportamento previsto por Albert, Jeong e Barabási (2000) para redes
\textit{scale-free}: tolerância a falhas aleatórias e fragilidade a ataques direcionados aos
hubs.

\subsection{Realocação de Demanda sob Capacidade}

A Tabela~\ref{tab:cascata} reporta a varredura de $\alpha$ no décimo passo do ataque, sobre a
\textbf{demanda total de 201.747.182 passageiros} (embarcados mais desembarcados --- o dobro do
volume de arestas, pois cada passageiro é contado na origem e no destino). No cenário de
referência $\alpha = 0{,}20$ (folga de 20\%), \textbf{115{,}9 milhões de passageiros (57{,}5\% da
demanda)} ficam sem realocação viável após o fechamento dos dez maiores hubs. A fração varia de
\textbf{29{,}1\% com folga de 100\%} a \textbf{61{,}0\% com folga de apenas 10\%}.

\begin{table}[h!]
\centering
\renewcommand{\arraystretch}{1.3}
\small
\begin{tabular}{cccc}
\toprule
\textbf{$\alpha$ (folga)} & \textbf{Aeroportos saturados} & \textbf{Demanda não realocada} & \textbf{\% da demanda} \\
\midrule
0{,}10 & 144 & 123.072.667 & 61{,}0\% \\
\rowcolor{azulclaro}
0{,}20 & 144 & 115.920.483 & 57{,}5\% \\
0{,}50 & 144 &  94.463.930 & 46{,}8\% \\
1{,}00 & 144 &  58.703.010 & 29{,}1\% \\
\bottomrule
\end{tabular}
\caption{Demanda não realocada no passo 10 do ataque, por folga de capacidade $\alpha$. A linha
destacada é o cenário de referência ($\alpha = 0{,}20$).}
\label{tab:cascata}
\end{table}

Dois achados, sempre condicionados às premissas do modelo (Seção~\ref{sec:limitacoes}):
(i)~a fração de demanda sem realocação é alta e cresce ao se reduzir a folga --- a remoção dos
grandes hubs introduz demanda em escala muito superior à capacidade residual do restante do
sistema; (ii)~o número de aeroportos saturados é praticamente constante ($\approx 144$) em toda
a faixa de $\alpha$ examinada, de modo que o sinal informativo é a \textbf{demanda não realocada},
não a contagem de saturados.

\paragraph{Resultado negativo que motivou o desenho.} A primeira versão do modelo seguiu a cascata
clássica por \textit{propagação de sobrecarga} (despejar a carga sobre vizinhos ignorando
capacidade, falhando quem excede o limite). O resultado é degenerado: a rede colapsa quase por
completo (resta um único aeroporto) e o desfecho é praticamente \textbf{insensível a $\alpha$} ---
dar 10\% ou 100\% de folga não muda o colapso. Esse comportamento, coerente com a observação de
Cumelles et al. (2021) de que modelos de fluxo contínuo são inadequados a redes aéreas, motivou a
substituição do mecanismo pela realocação que respeita a capacidade. A degenerescência é
reproduzível e está registrada como asserção de teste.

% 7.5 DISCUSSÃO
\section{Discussão: Contribuição, Inovação e Impacto}
\label{sec:discussao}

\subsection{Coerência dos Resultados}

Os três blocos de resultados são internamente coerentes e convergem com a literatura independente,
o que sustenta a validade da ferramenta. O ranking eleva ao topo aeroportos de forte papel de
intermediação --- com VCP em primeiro lugar, em concordância direta com Couto et al. (2015), que
já apontavam Viracopos como o nó cuja remoção mais fragmenta a malha nacional. A simulação
topológica reproduz a assimetria ataque/aleatório característica de redes \textit{scale-free}
(Albert et al., 2000). E o modelo de capacidade quantifica algo que a análise puramente
topológica não alcança: mesmo quando a rede permanece formalmente conectada, o sistema não dispõe
de capacidade residual para reabsorver a demanda dos hubs fechados. As três análises, obtidas por
caminhos independentes, contam a mesma história sobre a concentração de risco em poucos nós.

\subsection{Natureza da Contribuição: Inovação de Processo}

A contribuição central deste trabalho \textbf{não} é uma descoberta inédita sobre a topologia da
rede aérea brasileira --- essa estrutura já foi caracterizada academicamente. A inovação é de
\textbf{processo}: a entrega de um \textbf{artefato aberto, reproduzível e auditável} que organiza
microdados públicos da ANAC e métodos de análise de redes em uma ferramenta documentada e
executável. O panorama existente apresenta dois extremos --- de um lado, artigos acadêmicos que
\textit{descrevem} resultados mas não disponibilizam ferramenta executável e extensível; de outro,
soluções comerciais (OAG, Cirium) que são proprietárias e fechadas. Este trabalho ocupa a lacuna
intermediária: todo número deste relatório é regenerável a partir de \textit{scripts} versionados,
sem dependências externas (apenas a biblioteca padrão de Python; algoritmos de grafo escritos do
zero), acompanhado de visualização interativa offline. O valor está na \textbf{transformação de
dados dispersos em indicadores objetivos e reproduzíveis}.

Três elementos reforçam a solidez dessa contribuição, além da mera disponibilização de código:

\begin{itemize}[leftmargin=*]
  \item \textbf{Robustez metodológica.} O ranking de criticidade não depende de uma calibração
        arbitrária: VCP permanece como nó mais crítico sob todas as ponderações testadas e o top-10
        mantém-se estável aos pesos do índice (idêntico em cinco das seis perturbações; 9 dos 10
        retidos sob pesos uniformes). Conclusões que sobrevivem à variação de seus próprios
        parâmetros são mais confiáveis.
  \item \textbf{Um resultado negativo informativo.} A demonstração de que a cascata ingênua por
        propagação de sobrecarga colapsa a rede de forma insensível a $\alpha$ não é um fracasso
        --- é um achado que justifica empiricamente a escolha de um modelo de realocação com
        restrição de capacidade, e dialoga diretamente com Cumelles et al. (2021).
  \item \textbf{Honestidade epistêmica.} O trabalho distingue explicitamente \textbf{verificação}
        (consistência interna, garantida por testes automatizados) de \textbf{validação} (a
        correspondência com a realidade, declarada como trabalho futuro). Não se reivindica
        validação empírica que não foi feita --- o que torna as afirmações de confiabilidade
        precisas e defensáveis.
\end{itemize}

\subsection{Impacto}

No plano \textbf{prático}, o artefato oferece uma leitura objetiva de quais aeroportos concentram
risco estrutural na malha nacional, com subsídios para estudos de planejamento, contingência e
priorização de investimentos em capacidade. A camada de capacidade vai além da pergunta ``a rede
se desconecta?'' para ``o sistema consegue reabsorver os passageiros afetados?'' --- mais próxima
da experiência real do passageiro e do operador. No plano \textbf{científico e educacional}, por
ser aberto, testado e reproduzível, o artefato pode ser auditado, criticado e estendido por outros
pesquisadores e estudantes, servindo de base para os trabalhos futuros listados a seguir ---
notadamente a substituição do proxy de capacidade por dados reais. É precisamente a
reprodutibilidade que converte um exercício pontual em uma plataforma reutilizável.

% 7.6 LIMITAÇÕES E TRABALHOS FUTUROS
\section{Limitações e Trabalhos Futuros}
\label{sec:limitacoes}

A proposta delimita-se deliberadamente à análise da \textbf{vulnerabilidade estrutural} da rede, ou
seja, ao impacto da remoção de aeroportos sobre a \textit{topologia} da malha (conectividade,
componentes e isolamento de nós). Algumas escolhas de escopo merecem registro explícito:

\begin{itemize}[leftmargin=*]
  \item \textbf{Modelo estático.} O grafo representa a malha em um período de referência; não são
        modeladas dinâmicas temporais de voos (atrasos, propagação hora a hora), que exigiriam um
        modelo orientado a eventos, como discutido por Cumelles et al. (2021).
  \item \textbf{Capacidade como parâmetro, não como dado --- a premissa dos 20\%.} A limitação
        mais importante a registrar é o modelo de capacidade. Não existe base pública consolidada
        de capacidade real de pista e de \textit{slots} por aeroporto; por isso, a capacidade foi
        aproximada por $C_i = (1 + \alpha)\,L_i$, ou seja, supôs-se que cada aeroporto consegue
        absorver uma \textbf{folga de 20\% acima de sua própria demanda observada}
        ($\alpha = 0{,}20$) no cenário de referência. \textbf{Esse valor é uma suposição, não uma
        medida}: 20\% pode ser muito ou pouco, e variar amplamente entre um pequeno aeroporto
        regional e um grande hub metropolitano operando no limite. Justamente por reconhecer essa
        incerteza, os resultados não foram reportados para um único $\alpha$, mas ao longo de uma
        \textbf{varredura} (Tabela~\ref{tab:cascata}): a fração de demanda sem realocação varia de
        29\% (folga de 100\%) a 61\% (folga de 10\%). Essa sensibilidade tem duas leituras. A
        favor do trabalho: mesmo no cenário \textit{mais} otimista (100\% de folga), mais de um
        quarto da demanda fica desatendida --- o achado qualitativo (alta concentração de risco
        nos hubs) é \textbf{robusto à premissa}. Como limite: os \textbf{números absolutos} de
        demanda não realocada devem ser lidos como comparativos entre cenários, não como
        previsões. A \textbf{calibração empírica de $\alpha$} a partir de capacidades reais e de
        episódios históricos de fechamento de aeroportos é o trabalho futuro mais relevante, e a
        própria varredura indica por quê: é a faixa de $\alpha$ que determina a magnitude do
        impacto estimado.
  \item \textbf{Demais simplificações do modelo de capacidade.} A realocação opera sobre demanda
        anual agregada (não sobre voos programados, sem dimensão temporal) e é feita por
        proximidade geográfica, independentemente de o destino estar no mesmo componente conexo da
        rede fragmentada. A substituição da realocação por proximidade por um \textbf{modelo de
        Huff} com tempos de viagem (Correia et al.) permanece como refinamento futuro.
\end{itemize}

O reconhecimento explícito desses limites delimita com precisão o que o artefato entrega e o que
fica em aberto, orientando trabalhos subsequentes sem comprometer a validade dos resultados
obtidos.

% 8. CONCLUSÃO
\section{Conclusão}

Este trabalho entregou um artefato computacional aberto e reproduzível para a análise de
vulnerabilidade da rede aeroportuária brasileira, cumprindo os objetivos propostos. A partir de
microdados públicos da ANAC (doze meses de 2025), a malha foi modelada como um grafo ponderado de
155 aeroportos e 622 rotas, sobre o qual se calcularam um ranking de criticidade e simulações de
remoção de nós.

Os resultados respondem à questão orientadora. Os aeroportos estruturalmente mais críticos são os
de forte papel de intermediação --- com \textbf{Viracopos (VCP)} em primeiro lugar, em
concordância com a literatura independente --- e o impacto de sua remoção é fortemente assimétrico:
o ataque direcionado fragmenta a malha muito mais que a falha aleatória de mesma magnitude
(maior componente de 56{,}1\% contra 85{,}2\% no passo 10), confirmando a fragilidade típica de
redes \textit{scale-free}. A camada exploratória de capacidade acrescenta que, fechados os dez
maiores hubs, de 29\% a 61\% da demanda de passageiros fica sem realocação viável, conforme a
folga de capacidade assumida.

A contribuição é uma \textbf{inovação de processo}: não uma nova descoberta topológica, mas a
conversão de dados dispersos em uma ferramenta documentada, testada (verificação por testes
automatizados), robusta (ranking estável a perturbações nos pesos) e auditável. A principal fronteira em
aberto é a \textbf{validação empírica do modelo de capacidade} --- a substituição da folga
suposta de 20\% por capacidades reais e a calibração de $\alpha$ a partir de episódios históricos
de interrupção --- caminho que a natureza aberta e reproduzível do artefato deliberadamente
facilita.

% 9. REFERÊNCIAS
\section{Referências Bibliográficas}

\begin{itemize}[leftmargin=*, label={}, itemsep=8pt]

\item ALBERT, Reka; JEONG, Hawoong; BARABASI, Albert-Laszlo. \textbf{Error and attack tolerance
of complex networks.} \textit{Nature}, v. 406, p. 378--382, 2000. DOI: 10.1038/35019019.

\item CORREIA, Anderson Ribeiro; NIYAMA, Lucia Erika; NOGUEIRA, Sabrina de Araújo Furtado.
\textbf{Estimativa da distribuição da demanda na Região Metropolitana de São Paulo com cenários
de um novo aeroporto.} \textit{Journal of Transport Literature}, 2013.

\item COUTO, Guilherme S.; SILVA, Ana Paula Couto da; RUIZ, Linnyer B.; BENEVENUTO, Fabrício.
\textbf{Structural properties of the Brazilian air transportation network.} \textit{Anais da
Academia Brasileira de Ciências}, v. 87, n. 3, p. 1653--1674, 2015.
DOI: 10.1590/0001-3765201520140155.

\item CUMELLES, Joel; LORDAN, Oriol; SALLAN, Jose M. \textbf{Cascading failures in airport
networks.} \textit{Journal of Air Transport Management}, v. 92, 2021.
DOI: 10.1016/j.jairtraman.2021.102026.

\item AGÊNCIA NACIONAL DE AVIAÇÃO CIVIL (ANAC). \textbf{Dados de movimentação de passageiros
--- microdados de voos regulares.} Disponível em: \url{https://dados.anac.gov.br}. Acesso em:
abr. 2026.

\end{itemize}

\end{document}
